% Clase 7 — los cuatro primeros polinomios de Legendre en el intervalo que importa,
% con la normalizacion P_n(1)=1 que fija la convencion de la literatura.
\begin{tikzpicture}
\begin{axis}[emfig, width=10.5cm, height=6.6cm,
  xlabel={$u=\cos\theta$}, ylabel={$P_n(u)$},
  domain=-1:1, samples=120,
  xmin=-1.15, xmax=1.15, ymin=-1.35, ymax=1.55,
  xtick={-1,-0.5,0,0.5,1}, ytick={-1,-0.5,0,0.5,1},
  axis lines=middle,
  % OJO: `axis lines=middle` REDEFINE `every axis x label`, asi que el estilo
  % del rotulo tiene que ir DESPUES o queda centrado sobre la linea del eje
  % (y el `=` se lee como `\neq`).
  xlabel style={anchor=north west, yshift=-3pt},
  legend style={at={(0.5,-0.16)}, anchor=north, legend columns=4, draw=figgray!40}]
  \addplot[figblue]              {1};                    \addlegendentry{$P_0$}
  \addplot[figred]               {x};                    \addlegendentry{$P_1$}
  \addplot[figamber]             {0.5*(3*x^2-1)};        \addlegendentry{$P_2$}
  \addplot[figgray, dashed]      {0.5*(5*x^3-3*x)};      \addlegendentry{$P_3$}
  % la normalizacion: todos pasan por (1,1)
  \addplot[only marks, mark=*, mark size=1.6pt, figblue, forget plot]
    coordinates {(1,1)};
\end{axis}
\node[figlbl, align=center] at (5.2,-2.35)
  {Todos valen $1$ en $u=1$ (normalizaci\'on) y tienen la misma paridad que $n$.\\[0.2em]
   $P_n$ es de grado $n$: por eso son linealmente independientes.};
\end{tikzpicture}
